% Addition beside eq:alpha-recurrencia or inc-ampl-normalizacion.
% No formula is changed or removed, and no symbol is renamed.
\paragraph{Signed balance and sum with incoming quotient.}
\label{ampl-alfa-z-s-aclaracion}
It is useful to distinguish the two consecutive stages of positional
division by means of a common notation. The balance before quotient
transport is
\[
 Z_j:=P_j+E_j-\Phi_j-K_j^{\mathrm{per}}.
\]
The variable $s_j$ in \eqref{eq:alpha-recurrencia} denotes the sum that
already includes the incoming contribution, so that
\[
 s_j=Z_j+c_{j+1},\qquad
 A_j=Z_j+c_{j+1}-1000c_j.
\]
In \eqref{inc-ampl-normalizacion}, the letter $s_j$ locally denotes
the balance before that contribution, namely $Z_j$ in the present
notation. The two formulas therefore express the same Euclidean
division with differently named stages. The incidence configuration
$H^-_\alpha$ uses the signs of the balance $Z_j$ in the initial window;
the determination of the digits also uses the linking quotients
$c_{j+1}$ and $c_j$. This distinction jointly preserves the signed
support and the normalization of the expansion.
